\section{Contrastive alignment and what it identifies}
\label{sec:C05:contrastive}
CLIP learns image and text representations through a contrastive objective on paired data~\cite{radford2021clip}. Let normalized embeddings be $u_i$ and $v_j$ and temperature $\tau>0$. A directional batch loss is
\begin{equation}
 -\frac1B\sum_{i=1}^{B}\log
 \frac{\exp(u_i^\top v_i/\tau)}
 {\sum_{j=1}^{B}\exp(u_i^\top v_j/\tau)}.
\end{equation}
The symmetric objective also predicts images from texts. This trains compatibility between paired modalities relative to the negative set. It does not by itself supervise a state transition, a tool argument, or the consequences of an action. A representation that matches a screenshot to a caption can still confuse whether a visible operation has already occurred.

Negative construction affects the learned distinctions. Two captions can both accurately describe an image while the batch treats only one as its paired positive. Conversely, dataset artifacts can make a pair identifiable without learning the causal structure relevant to control. These issues do not invalidate contrastive learning; they identify its objective. An agent evaluation should therefore test exact distinctions such as enabled versus disabled, before versus after, or selected versus merely visible. Generic image-text retrieval can be included as a diagnostic, but should not substitute for the closed-loop decision task.

Temperature controls the sharpness of the softmax over similarities. A smaller temperature concentrates gradients on hard alternatives and can magnify calibration differences between domains. The similarity score is not a calibrated probability that a proposed tool action is correct. Thresholding it as though it were one can yield unstable rejection behavior after a change in image style, prompt wording, or candidate set. Calibration should be measured against an explicit event, such as whether a retrieved frame contains the target control, under the deployment candidate distribution.

Cross-modal alignment also has a temporal dimension. A transcript can describe an earlier state while an image shows the current one. Combining them without timestamps can produce a fluent but internally inconsistent account. The representation pipeline should preserve modality, source, time, and transformation lineage. If an OCR engine and a vision-language model disagree, the agent needs a policy for seeking additional evidence rather than silently trusting whichever channel produces more confident language. Such disagreement is often an informative signal that the observation boundary needs another interaction.

